\chapter*{Notation}
\addcontentsline{toc}{chapter}{Notation}
\markboth{Notation}{Notation}

All numbers are relative to the \textbf{key} of the tune: 1 is the tonic.
The same tune can be studied in any key; the \textbf{Keys} line lists the
keys commonly used (original key first).

\section*{Melody}
\renewcommand{\arraystretch}{1.5}
\begin{tabular}{p{2.4cm}p{2.8cm}p{6cm}}
\textbf{Symbol} & \textbf{Example} & \textbf{Meaning} \\ \hline
\mel{1 \ldots 7} & \mel{1\;2\;3} & Scale degrees of the key (in F: F G A). \\
\mel{\flat n}, \mel{\sharp n} & \mel{\flat3\;\sharp4} & Altered degrees (in F: A$\flat$, B$\natural$). \\
\mel{\hi{n}} & \mel{\hi{1}} & Dot above: one octave higher. \\
\mel{\lo{n}} & \mel{\lo{7}} & Dot below: one octave lower. \\
\mel{n} (no dot) & \mel{5} & Middle (reference) octave. \\
\tg{n} & \mel{\tg{3}} & Circled: target / landing note (e.g.\ guide tone). \\
\mel{\tg{n}}\tfn{f} & \mel{\tg{\hi{1}}}\tfn{$\sharp$9} & Circled note over a chromatic chord, with its function on that chord (in red). \\
\mel{-} & \mel{3\;-} & Previous note held. \\
\end{tabular}

\section*{Tune}
The \textbf{Tune} section gives the melody with its rhythm: scale degrees on a single
line, with the stems, beams, flags, dots, ties and rests of a score (half note: degree in
a rounded box). Eighth notes are swung unless stated. Accents: \textasciicircum\ short,
$>$ accented. Chords are placed on the beat where they start; grey: alternative changes.

\section*{Chords}
\begin{tabular}{p{2.4cm}p{2.8cm}p{6cm}}
\textbf{Symbol} & \textbf{Example (F)} & \textbf{Meaning} \\ \hline
Roman numeral & \ch{IV}{\Delta7} & Chord built on that degree of the key (here B$\flat$). \\
\ch{\flat VII}{7} & E$\flat$7 & Chord on an altered degree. \\
\ch{}{\Delta7} & \ch{I}{\Delta7} = F$\Delta$7 & Major seventh. \\
\ch{}{-7} & \ch{ii}{-7} = G$-$7 & Minor seventh. \\
\ch{}{7} & \ch{V}{7} = C7 & Dominant seventh. \\
\ch{}{\text{\o}7} & \ch{vii}{\text{\o}7} = E\o7 & Half-diminished (m7$\flat$5). \\
\ch{}{\circ7} & \ch{\sharp i}{\circ7} = F$\sharp^{\circ}$7 & Diminished seventh. \\
\chx{VI}{7} & \chx{VI}{7} = G7 & Chord outside the key (secondary dominant, borrowed chord, \ldots), in red. \\
\alt{\ch{VI}{7}} & \alt{\ch{V}{7}/\ch{ii}{}} & Alternative changes (secondary dominants, substitutions): smaller, grey, below the main changes. \\
\end{tabular}

\medskip
Upper case numerals for major and dominant chords, lower case for minor,
half-diminished and diminished chords.

\section*{Chord functions of melody notes}
\begin{tabular}{p{2.4cm}p{2.8cm}p{6cm}}
\textbf{Symbol} & \textbf{Example} & \textbf{Meaning} \\ \hline
\fn{R 3 5 7} & \fn{R $\flat$3 5 $\flat$7} & Interval of each melody note above the
\emph{chord} root (not the key). Extensions: \fn{9 11 13}, altered:
\fn{$\flat$9 $\sharp$11 $\flat$13}. \\
\end{tabular}

\section*{Form grid}
\begin{form}
A & \mel{2\;4\;6\;\tg{\hi{1}}} & \mel{5\;\tg{7}\;\hi{2}\;\hi{4}} & \mel{\tg{\hi{3}}} & \mel{-} \\
  & \ch{ii}{-7} & \ch{V}{7} & \ch{I}{\Delta7} & \ch{I}{\Delta7} \\
  & & & & \alt{\ch{I}{\Delta7}\; \ch{VI}{7}} \\
  & \fn{R $\flat$3 5 $\flat$7} & \fn{R 3 5 $\flat$7} & \fn{3} & \\ \hline
\end{form}

\medskip
\begin{itemize}
\item Each cell is one bar; four bars per row.
\item Left column: section label (A, B, \ldots).
\item Rows: melody (scale degrees, top); chords (bold); alternative changes (grey, optional); chord function of the melody notes.
\end{itemize}

\section*{Boxes}
\begin{study}{Study box}
What to study in a passage: patterns, guide tones, licks, exercises.
\end{study}

\begin{notes}
General notes: history of the tune, composers, context, anecdotes.
\end{notes}

\begin{patternbox}{Name of the pattern}{chords it applies to}
A short phrase for improvising, in scale degrees of the key of the target chord.
Patterns are shared between tunes; each tune lists the bars where they apply.
\end{patternbox}

\begin{recordings}
\rec{Artist}{Album}{year} -- recordings worth listening to and transcribing.
\end{recordings}
